% GENERATED FILE. Edit canonical lessons, then run npm run generate:books.
% canonical-source-hash: fnv1a64:175308d0165003ce
% canonical-lessons: JA-C458-bengoshi, JA-C458-koumuin, JA-C458-haiyuu, JA-C458-sakka, JA-C458-shokunin

\chapter{Jobs People Do}
\label{ch:ja-jobs-people-do}
\hlchaptermodality{458}

\begin{chapteropening}
\textbf{By the end of this chapter:} \emph{I can name some jobs people do.}

Complete the last lesson of chapter 458: I can name some jobs people do.
\end{chapteropening}

\section[\texorpdfstring{\ja{べんごし} (bengoshi)}{べんごし (bengoshi)}]{\ja{べんごし} (bengoshi) — a lawyer}

\label{lesson:JA-C458-bengoshi}



\begin{quote}
Before the new one: say the Japanese for to apply, then the Japanese for recruitment.
\end{quote}

\subsection*{You'll want to know: \ja{べんごし}}
\textbf{\ja{べんごし}} — \emph{bengoshi} — \textquotedblleft{}a lawyer\textquotedblright{}.

\begin{quote}
\textbf{\ja{あには} \ja{べんごしに} \ja{なりたいです。}} — \emph{ani wa bengoshi ni naritai desu} — \textquotedblleft{}My elder brother wants to become a lawyer.\textquotedblright{}
\end{quote}

\begin{cousinweb}
\textbf{\ja{べんごし}} (\emph{bengoshi}) is Sino-Japanese (kango): three Chinese characters read with their on'yomi, readings borrowed from Chinese. They are \emph{ben}, \textquotedblleft{}speech\textquotedblright{} (Middle Chinese \emph{bjenX}), \emph{go}, \textquotedblleft{}protect\textquotedblright{} (\emph{huH}), and \emph{shi}, \textquotedblleft{}gentleman\textquotedblright{} (\emph{dzriX}) (KANJIDIC2; Middle Chinese in Baxter's transcription).
\end{cousinweb}

\begin{grammarlens}[title={What you've built}]
One more word for describing what you lived and felt.
\end{grammarlens}

\subsection*{Your turn}
\begin{itemize}
\raggedright
  \item \emph{Say it:} \emph{bengoshi}
  \item \emph{Say it:} \emph{bengoshi}, once more
  \item \emph{Say it:} \emph{ani wa bengoshi ni naritai desu} — My elder brother wants to become a lawyer.
  \item \emph{Recall:} say the Japanese for to apply, then the Japanese for recruitment, then say \emph{bengoshi} again
\end{itemize}

\begin{tcolorbox}[breakable,colback=teal!4,colframe=teal!35!black,title={Before you move on}]
What does \ja{べんごし} mean? (\textquotedblleft{}A lawyer\textquotedblright{}.) Say it once more.
\end{tcolorbox}
\section[\texorpdfstring{\ja{こうむいん} (kōmuin)}{こうむいん (kōmuin)}]{\ja{こうむいん} (kōmuin) — a civil servant}

\label{lesson:JA-C458-koumuin}



\begin{quote}
Before the new one: say the Japanese for recruitment, then the Japanese for a lawyer.
\end{quote}

\subsection*{You'll want to know: \ja{こうむいん}}
\textbf{\ja{こうむいん}} — \emph{kōmuin} — \textquotedblleft{}a civil servant\textquotedblright{}.

\begin{quote}
\textbf{\ja{ちちは} \ja{しやくしょの} \ja{こうむいんです。}} — \emph{chichi wa shiyakusho no kōmuin desu} — \textquotedblleft{}My father is a civil servant at the city hall.\textquotedblright{}
\end{quote}

\begin{cousinweb}
\textbf{\ja{こうむいん}} (\emph{kōmuin}) is Sino-Japanese (kango): three Chinese characters read with their on'yomi, readings borrowed from Chinese. They are \emph{kō}, \textquotedblleft{}public\textquotedblright{} (Middle Chinese \emph{kuwng}), \emph{mu}, \textquotedblleft{}duties\textquotedblright{} (\emph{mjuH}), and \emph{in}, \textquotedblleft{}member\textquotedblright{} (\emph{hjwen}) (KANJIDIC2; Middle Chinese in Baxter's transcription).
\end{cousinweb}

\begin{grammarlens}[title={What you've built}]
One more word for describing what you lived and felt.
\end{grammarlens}

\subsection*{Your turn}
\begin{itemize}
\raggedright
  \item \emph{Say it:} \emph{kōmuin}
  \item \emph{Say it:} \emph{kōmuin}, once more
  \item \emph{Say it:} \emph{chichi wa shiyakusho no kōmuin desu} — My father is a civil servant at the city hall.
  \item \emph{Recall:} say the Japanese for recruitment, then the Japanese for a lawyer, then say \emph{kōmuin} again
\end{itemize}

\begin{tcolorbox}[breakable,colback=teal!4,colframe=teal!35!black,title={Before you move on}]
What does \ja{こうむいん} mean? (\textquotedblleft{}A civil servant\textquotedblright{}.) Say it once more.
\end{tcolorbox}
\section[\texorpdfstring{\ja{はいゆう} (haiyū)}{はいゆう (haiyū)}]{\ja{はいゆう} (haiyū) — an actor}

\label{lesson:JA-C458-haiyuu}



\begin{quote}
Before the new one: say the Japanese for a lawyer, then the Japanese for a civil servant.
\end{quote}

\subsection*{You'll want to know: \ja{はいゆう}}
\textbf{\ja{はいゆう}} — \emph{haiyū} — \textquotedblleft{}an actor\textquotedblright{}.

\begin{quote}
\textbf{\ja{あの} \ja{はいゆうの} \ja{えいがを} \ja{みました。}} — \emph{ano haiyū no eiga o mimashita} — \textquotedblleft{}I saw a film with that actor.\textquotedblright{}
\end{quote}

\begin{cousinweb}
\textbf{\ja{はいゆう}} (\emph{haiyū}) is Sino-Japanese (kango): two Chinese characters read with their on'yomi, readings borrowed from Chinese. They are \emph{hai}, \textquotedblleft{}actor\textquotedblright{} (Middle Chinese \emph{beaj}) and \emph{yū}, \textquotedblleft{}actor\textquotedblright{} (\emph{'juw}) (KANJIDIC2; Middle Chinese in Baxter's transcription).
\end{cousinweb}

\begin{grammarlens}[title={What you've built}]
One more word for describing what you lived and felt.
\end{grammarlens}

\subsection*{Your turn}
\begin{itemize}
\raggedright
  \item \emph{Say it:} \emph{haiyū}
  \item \emph{Say it:} \emph{haiyū}, once more
  \item \emph{Say it:} \emph{ano haiyū no eiga o mimashita} — I saw a film with that actor.
  \item \emph{Recall:} say the Japanese for a lawyer, then the Japanese for a civil servant, then say \emph{haiyū} again
\end{itemize}

\begin{tcolorbox}[breakable,colback=teal!4,colframe=teal!35!black,title={Before you move on}]
What does \ja{はいゆう} mean? (\textquotedblleft{}An actor\textquotedblright{}.) Say it once more.
\end{tcolorbox}
\section[\texorpdfstring{\ja{さっか} (sakka)}{さっか (sakka)}]{\ja{さっか} (sakka) — a novelist; an author}

\label{lesson:JA-C458-sakka}



\begin{quote}
Before the new one: say the Japanese for a civil servant, then the Japanese for an actor.
\end{quote}

\subsection*{You'll want to know: \ja{さっか}}
\textbf{\ja{さっか}} — \emph{sakka} — \textquotedblleft{}a novelist; an author\textquotedblright{}.

\begin{quote}
\textbf{\ja{ははは} \ja{あの} \ja{さっかが} \ja{すきです。}} — \emph{haha wa ano sakka ga suki desu} — \textquotedblleft{}My mother likes that author.\textquotedblright{}
\end{quote}

\begin{cousinweb}
\textbf{\ja{さっか}} (\emph{sakka}) is Sino-Japanese (kango): two Chinese characters read with their on'yomi, readings borrowed from Chinese. They are \emph{sa}, \textquotedblleft{}make\textquotedblright{} (Middle Chinese \emph{tsak}) and \emph{ka}, \textquotedblleft{}expert\textquotedblright{} (\emph{kae}) (KANJIDIC2; Middle Chinese in Baxter's transcription).
\end{cousinweb}

\begin{grammarlens}[title={What you've built}]
One more word for describing what you lived and felt.
\end{grammarlens}

\subsection*{Your turn}
\begin{itemize}
\raggedright
  \item \emph{Say it:} \emph{sakka}
  \item \emph{Say it:} \emph{sakka}, once more
  \item \emph{Say it:} \emph{haha wa ano sakka ga suki desu} — My mother likes that author.
  \item \emph{Recall:} say the Japanese for a civil servant, then the Japanese for an actor, then say \emph{sakka} again
\end{itemize}

\begin{tcolorbox}[breakable,colback=teal!4,colframe=teal!35!black,title={Before you move on}]
What does \ja{さっか} mean? (\textquotedblleft{}A novelist; an author\textquotedblright{}.) Say it once more.
\end{tcolorbox}
\section[\texorpdfstring{\ja{しょくにん} (shokunin)}{しょくにん (shokunin)}]{\ja{しょくにん} (shokunin) — a craftsman; an artisan}

\label{lesson:JA-C458-shokunin}



\begin{quote}
Before the new one: say the Japanese for an actor, then the Japanese for a novelist.
\end{quote}

\subsection*{You'll want to know: \ja{しょくにん}}
\textbf{\ja{しょくにん}} — \emph{shokunin} — \textquotedblleft{}a craftsman; an artisan\textquotedblright{}.

\begin{quote}
\textbf{\ja{そふは} \ja{たたみの} \ja{しょくにんでした。}} — \emph{sofu wa tatami no shokunin deshita} — \textquotedblleft{}My grandfather was a tatami craftsman.\textquotedblright{}
\end{quote}

\begin{cousinweb}
\textbf{\ja{しょくにん}} (\emph{shokunin}) is Sino-Japanese (kango): two Chinese characters read with their on'yomi, readings borrowed from Chinese. They are \emph{shoku}, \textquotedblleft{}employment\textquotedblright{} (Middle Chinese \emph{tsyik}) and \emph{nin}, \textquotedblleft{}person\textquotedblright{} (\emph{nyin}) (KANJIDIC2; Middle Chinese in Baxter's transcription).
\end{cousinweb}

\begin{grammarlens}[title={What you've built}]
One more word for describing what you lived and felt.
\end{grammarlens}

\subsection*{Your turn}
\begin{itemize}
\raggedright
  \item \emph{Say it:} \emph{shokunin}
  \item \emph{Say it:} \emph{shokunin}, once more
  \item \emph{Say it:} \emph{sofu wa tatami no shokunin deshita} — My grandfather was a tatami craftsman.
  \item \emph{Recall:} say the Japanese for an actor, then the Japanese for a novelist, then say \emph{shokunin} again
\end{itemize}

\begin{tcolorbox}[breakable,colback=teal!4,colframe=teal!35!black,title={Before you move on}]
What does \ja{しょくにん} mean? (\textquotedblleft{}A craftsman; an artisan\textquotedblright{}.) Say it once more.
\end{tcolorbox}
